% GENERATED FILE. Edit canonical lessons, then run npm run generate:books.
% canonical-source-hash: fnv1a64:7c18831a53a97c98
% canonical-lessons: PA-C230-mokla, PA-C230-lahuna, PA-C230-lapetna, PA-C230-sutthan, PA-C230-jhumka

\chapter{Getting Dressed}
\label{ch:pa-a2-getting-dressed}
\hlchaptermodality{230}

\begin{chapteropening}
\textbf{By the end of this chapter:} \emph{I can talk about putting on and taking off clothes.}

Complete the last lesson of chapter 230: I can talk about putting on and taking off clothes.
\end{chapteropening}

\section[\texorpdfstring{\pa{ਮੋਕਲਾ} (moklā)}{ਮੋਕਲਾ (moklā)}]{\pa{ਮੋਕਲਾ} (moklā) — roomy; loose-fitting}

\label{lesson:PA-C230-mokla}



\begin{quote}
Before the new one: say the Punjabi for a cord, then the Punjabi for a knotted cloth button.
\end{quote}

\subsection*{You'll want to know: \pa{ਮੋਕਲਾ}}
\textbf{\pa{ਮੋਕਲਾ}} — \emph{moklā} — \textquotedblleft{}roomy; loose-fitting\textquotedblright{}.

\begin{quote}
\pa{ਇਹ} \pa{ਕੁੜਤਾ} \pa{ਮੋਕਲਾ} \pa{ਹੈ} \pa{ਪਰ} \pa{ਉਹ} \pa{ਕਮੀਜ਼} \pa{ਭੀੜੀ} \pa{ਹੈ।}
\end{quote}

\begin{quote}
\emph{eh kuṛtā moklā hai par oh kamīz bhīṛī hai.}
\end{quote}

\begin{quote}
This kurta is roomy, but that shirt is tight.
\end{quote}

\begin{cousinweb}
\textbf{\pa{ਮੋਕਲਾ}} (\emph{moklā}) is inherited from a reconstructed \emph{*mukna}, \textquotedblleft{}loosed\textquotedblright{}; CDIAL lists Punjabi \emph{moklā}, \textquotedblleft{}loose\textquotedblright{} (CDIAL 10157).
\end{cousinweb}

\begin{grammarlens}[title={What you've built}]
One more word for saying what you do.
\end{grammarlens}

\subsection*{Your turn}
\begin{itemize}
\raggedright
  \item \emph{Say it:} \emph{moklā}
  \item \emph{Say it:} \emph{moklā}, once more
  \item \emph{Say it:} \emph{eh kuṛtā moklā hai par oh kamīz bhīṛī hai.} — This kurta is roomy, but that shirt is tight.
  \item \emph{Recall:} say the Punjabi for a cord, then the Punjabi for a knotted cloth button, then say \emph{moklā} again
\end{itemize}

\begin{tcolorbox}[breakable,colback=teal!4,colframe=teal!35!black,title={Before you move on}]
What does \pa{ਮੋਕਲਾ} mean? (\textquotedblleft{}Roomy; loose-fitting\textquotedblright{}.) Say it once more.
\end{tcolorbox}
\section[\texorpdfstring{\pa{ਲਾਹੁਣਾ} (lāhuṇā)}{ਲਾਹੁਣਾ (lāhuṇā)}]{\pa{ਲਾਹੁਣਾ} (lāhuṇā) — to take off (clothes or shoes)}

\label{lesson:PA-C230-lahuna}



\begin{quote}
Before the new one: say the Punjabi for a knotted cloth button, then the Punjabi for roomy.
\end{quote}

\subsection*{You'll want to know: \pa{ਲਾਹੁਣਾ}}
\textbf{\pa{ਲਾਹੁਣਾ}} — \emph{lāhuṇā} — \textquotedblleft{}to take off (clothes or shoes)\textquotedblright{}.

\begin{quote}
\pa{ਮੈਂ} \pa{ਘਰ} \pa{ਵਿੱਚ} \pa{ਜੁੱਤੀ} \pa{ਲਾਹੁੰਦਾ} \pa{ਹਾਂ।}
\end{quote}

\begin{quote}
\emph{maiṁ ghar vich juttī lāhundā hāṁ.}
\end{quote}

\begin{quote}
I take off my shoes in the house.
\end{quote}

Past (I …): \pa{ਮੈਂ} \pa{ਕੋਟ} \pa{ਲਾਹਿਆ}, \emph{maiṁ koṭ lāhiā}.

\begin{cousinweb}
\textbf{\pa{ਲਾਹੁਣਾ}} (\emph{lāhuṇā}) is inherited from a reconstructed \emph{*lāsayati}, \textquotedblleft{}makes slip, makes descend\textquotedblright{}; CDIAL lists Punjabi \emph{lāhṇā}, \textquotedblleft{}to bring down, unload\textquotedblright{} (CDIAL 11042), but not the everyday sense \textquotedblleft{}take off (clothes)\textquotedblright{}.
\end{cousinweb}

\begin{grammarlens}[title={What you've built}]
One more word for saying what you do.
\end{grammarlens}

\subsection*{Your turn}
\begin{itemize}
\raggedright
  \item \emph{Say it:} \emph{lāhuṇā}
  \item \emph{Say it:} \emph{lāhuṇā}, once more
  \item \emph{Say it:} \emph{maiṁ ghar vich juttī lāhundā hāṁ.} — I take off my shoes in the house.
  \item \emph{Recall:} say the Punjabi for a knotted cloth button, then the Punjabi for roomy, then say \emph{lāhuṇā} again
  \item \emph{Write it:} copy \textbf{\pa{ਲਾਹੁਣਾ}} once, checking the whole word
\end{itemize}

\begin{tcolorbox}[breakable,colback=teal!4,colframe=teal!35!black,title={Before you move on}]
What does \pa{ਲਾਹੁਣਾ} mean? (\textquotedblleft{}To take off (clothes or shoes)\textquotedblright{}.) Say it once more.
\end{tcolorbox}
\section[\texorpdfstring{\pa{ਲਪੇਟਣਾ} (lapeṭṇā)}{ਲਪੇਟਣਾ (lapeṭṇā)}]{\pa{ਲਪੇਟਣਾ} (lapeṭṇā) — to wrap}

\label{lesson:PA-C230-lapetna}



\begin{quote}
Before the new one: say the Punjabi for roomy, then the Punjabi for to take off.
\end{quote}

\subsection*{You'll want to know: \pa{ਲਪੇਟਣਾ}}
\textbf{\pa{ਲਪੇਟਣਾ}} — \emph{lapeṭṇā} — \textquotedblleft{}to wrap\textquotedblright{}.

\begin{quote}
\pa{ਸਰਦੀ} \pa{ਵਿੱਚ} \pa{ਮੈਂ} \pa{ਕੰਬਲ} \pa{ਲਪੇਟਦੀ} \pa{ਹਾਂ।}
\end{quote}

\begin{quote}
\emph{sardī vich maiṁ kambal lapeṭdī hāṁ.}
\end{quote}

\begin{quote}
In winter I wrap myself in a blanket.
\end{quote}

Past (I …): \pa{ਮੈਂ} \pa{ਚੁੰਨੀ} \pa{ਲਪੇਟੀ}, \emph{maiṁ chunnī lapeṭī}.

\begin{cousinweb}
\textbf{\pa{ਲਪੇਟਣਾ}} (\emph{lapeṭṇā}) is inherited from a reconstructed \emph{*lappēṭṭ}, \textquotedblleft{}wrap\textquotedblright{}; CDIAL lists Punjabi \emph{lapeṭnā}, \textquotedblleft{}to wrap\textquotedblright{} (CDIAL 10942).
\end{cousinweb}

\begin{grammarlens}[title={What you've built}]
One more word for saying what you do.
\end{grammarlens}

\subsection*{Your turn}
\begin{itemize}
\raggedright
  \item \emph{Say it:} \emph{lapeṭṇā}
  \item \emph{Say it:} \emph{lapeṭṇā}, once more
  \item \emph{Say it:} \emph{sardī vich maiṁ kambal lapeṭdī hāṁ.} — In winter I wrap myself in a blanket.
  \item \emph{Recall:} say the Punjabi for roomy, then the Punjabi for to take off, then say \emph{lapeṭṇā} again
  \item \emph{Write it:} copy \textbf{\pa{ਲਪੇਟਣਾ}} once, checking the whole word
\end{itemize}

\begin{tcolorbox}[breakable,colback=teal!4,colframe=teal!35!black,title={Before you move on}]
What does \pa{ਲਪੇਟਣਾ} mean? (\textquotedblleft{}To wrap\textquotedblright{}.) Say it once more.
\end{tcolorbox}
\section[\texorpdfstring{\pa{ਸੁੱਥਣ} (sutthaṇ)}{ਸੁੱਥਣ (sutthaṇ)}]{\pa{ਸੁੱਥਣ} (sutthaṇ) — loose trousers (worn by women)}

\label{lesson:PA-C230-sutthan}



\begin{quote}
Before the new one: say the Punjabi for to take off, then the Punjabi for to wrap.
\end{quote}

\subsection*{You'll want to know: \pa{ਸੁੱਥਣ}}
\textbf{\pa{ਸੁੱਥਣ}} — \emph{sutthaṇ} — \textquotedblleft{}loose trousers (worn by women)\textquotedblright{}.

\begin{quote}
\pa{ਦਾਦੀ} \pa{ਸੁੱਥਣ} \pa{ਅਤੇ} \pa{ਕੁੜਤਾ} \pa{ਪਹਿਨਦੀ} \pa{ਹੈ।}
\end{quote}

\begin{quote}
\emph{dādī sutthaṇ ate kuṛtā pahindī hai.}
\end{quote}

\begin{quote}
Grandmother wears loose trousers and a kurta.
\end{quote}

A feminine noun.

\begin{cousinweb}
\textbf{\pa{ਸੁੱਥਣ}} (\emph{sutthaṇ}) is inherited from a reconstructed \emph{*sutthanā}, \textquotedblleft{}trousers\textquotedblright{}; CDIAL lists Punjabi \emph{sutthaṇ}, \textquotedblleft{}women's loose trousers\textquotedblright{} (CDIAL 13468).
\end{cousinweb}

\begin{grammarlens}[title={What you've built}]
One more word for saying what you do.
\end{grammarlens}

\subsection*{Your turn}
\begin{itemize}
\raggedright
  \item \emph{Say it:} \emph{sutthaṇ}
  \item \emph{Say it:} \emph{sutthaṇ}, once more
  \item \emph{Say it:} \emph{dādī sutthaṇ ate kuṛtā pahindī hai.} — Grandmother wears loose trousers and a kurta.
  \item \emph{Recall:} say the Punjabi for to take off, then the Punjabi for to wrap, then say \emph{sutthaṇ} again
  \item \emph{Write it:} copy \textbf{\pa{ਸੁੱਥਣ}} once, checking the whole word
\end{itemize}

\begin{tcolorbox}[breakable,colback=teal!4,colframe=teal!35!black,title={Before you move on}]
What does \pa{ਸੁੱਥਣ} mean? (\textquotedblleft{}Loose trousers (worn by women)\textquotedblright{}.) Say it once more.
\end{tcolorbox}
\section[\texorpdfstring{\pa{ਝੁਮਕਾ} (jhumkā)}{ਝੁਮਕਾ (jhumkā)}]{\pa{ਝੁਮਕਾ} (jhumkā) — a hanging earring}

\label{lesson:PA-C230-jhumka}



\begin{quote}
Before the new one: say the Punjabi for to wrap, then the Punjabi for loose trousers.
\end{quote}

\subsection*{You'll want to know: \pa{ਝੁਮਕਾ}}
\textbf{\pa{ਝੁਮਕਾ}} — \emph{jhumkā} — \textquotedblleft{}a hanging earring\textquotedblright{}.

\begin{quote}
\pa{ਭੈਣ} \pa{ਦੇ} \pa{ਝੁਮਕੇ} \pa{ਲਾਲ} \pa{ਹਨ।}
\end{quote}

\begin{quote}
\emph{bhaiṇ de jhumke lāl han.}
\end{quote}

\begin{quote}
My sister's earrings are red.
\end{quote}

A masculine noun.

\begin{cousinweb}
\textbf{\pa{ਝੁਮਕਾ}} (\emph{jhumkā}) is inherited from a reconstructed \emph{*jhumma}, a form of \emph{*jhuppa}, \textquotedblleft{}crowd, cluster\textquotedblright{}; CDIAL lists Punjabi \emph{jhumkā}, \textquotedblleft{}pendent earring\textquotedblright{} (CDIAL 5404).
\end{cousinweb}

\begin{grammarlens}[title={What you've built}]
One more word for saying what you do.
\end{grammarlens}

\subsection*{Your turn}
\begin{itemize}
\raggedright
  \item \emph{Say it:} \emph{jhumkā}
  \item \emph{Say it:} \emph{jhumkā}, once more
  \item \emph{Say it:} \emph{bhaiṇ de jhumke lāl han.} — My sister's earrings are red.
  \item \emph{Recall:} say the Punjabi for to wrap, then the Punjabi for loose trousers, then say \emph{jhumkā} again
  \item \emph{Write it:} copy \textbf{\pa{ਝੁਮਕਾ}} once, checking the whole word
\end{itemize}

\begin{tcolorbox}[breakable,colback=teal!4,colframe=teal!35!black,title={Before you move on}]
What does \pa{ਝੁਮਕਾ} mean? (\textquotedblleft{}A hanging earring\textquotedblright{}.) Say it once more.
\end{tcolorbox}
